 \section{Motivation, goal of the project, research question}
    \chapterauthor{Agnieszka Kuleta, Kamil Dędza}
    
    Food is an intrinsic part of our lives, without it, we simply cannot survive. We eat not only to provide our bodies with the essential nutrients they need, but also because food brings us some kind of pleasure.  But what if your body reacts badly to certain substances? The answer is straightforward: we can stop consuming them. This approach works well if you need to avoid just one or two ingredients, we just eliminate certain meals from our diet. But what happens if you have to give up several of your favorite ingredients and you do not know how to substitute them so the taste of your favorite meals remains unchanged? Would you completely change your diets without any substitutes of your favorite dishes? Again everyone will answer that they would start to look for some substitutes or completely new recipes. However, finding them is not that easy.
    {\par}
    Nowadays, the number of recipes available for people with food restrictions is quite limited and the number of people with many specific food allergies \cite{warren2020epidemiology} and specific diets is still growing as shown on the Figure \ref{fig:special_diets}, \cite{martin2020prescription}.
    
    \begin{figure}[ht]
    \centering
    \includegraphics[width=0.8\textwidth]{special_diets_graph.png}
    \caption{Trends in the use of special diets among U.S. adults, 2007--2018 \cite{martin2020prescription}.}
    \label{fig:special_diets}
\end{figure}
    People, who are forced to follow the special diet, often eat only what they can, not what they want, because finding suitable substitutes is challenging,  \cite{foodiversity2024} without going through thousands of books or blogs. These people are just tired of trying new things. This lack of accessible and varied recipes limits their dietary options, making it difficult to enjoy eating and could cause serious health issues  \cite{foodiversity2024}. Here raises the question: How can we add variety to our diet and eat like everyone else without showing any signs of helplessness, dissatisfaction or serious health problems?
    {\par}
    As you can clearly see that there is a growing need for more resources and creative solutions to the culinary world. Our culinary portal “Foodies” is created to fill this gap. Our aim was to create a  web application, which is not only a place for suggestions of the next breakfast or dinner, but also a help in managing your specific diets and everyone’s needs by eliminating all unwanted ingredients and substituting them with one which you can and like to eat.
    {\par}
    Existing solutions available on the market are mainly concentrated on the recipes itself, with little to no help for people with more specific needs. As you can see in further research \ref{section:existing_solutions}, none of these solutions provides directly different substitutions of ingredients. Some portals write the substitute for one ingredient, mostly to substitute meat’s product, but the substitute is not included in the number of calories or other nutrition. Furthermore, changing  wheat  or dairy products was not the option. To find some substitutes, they need to search many different sources, which is time consuming and inconvenient. Moreover, not even all portals provided the numbers of calories or nutrition in the recipe \ref{section:existing_solutions}.People want to eat what they can and like and want to know how much they eat.\cite{howweeat} So, the goal of our project is to enable them not to give up what they want to eat, only because some of the ingredients do not match theirs’ diet. 
    {\par}
    For many years, our whole team has been interested in nutrition, each of us with a different perspective.  Cooking, calories intake, meal preparation to name a few, are topics that our group is familiar with. Combined, we form a solid team to put our expertise for a good use. The project will focus on the users with specific diets, which are becoming more and more popular nowadays. The app will let them personalize their searches with the type of the diet they are currently following, rate and comment these recipes and personalize their profiles.
        
    \section{Project organization}
    \chapterauthor{Agnieszka Kuleta}
    
    Having said about our motivation to do this project, let's talk about implementation.The implementation of the project is divided in three key stages: business analysis, implementation, test and verification. 
    \par
    The business analysis stage started from setting our goal and our motivation of doing this project, as was outlined in this chapter.  We gathered all requirements, which will be detailed in chapter \ref{chapter:requirements}. Additionally, we chose a project development methodology, unanimously deciding on \textbf{Scrum}, due to its iterative and collaborative framework. After establishing tools, team responsibilities and initiating our Backlog, which we will describe more comprehensively in the chapter \ref{chapter:methodology}.
    \par
    Completing business stage brought us to implementation stage, which began with data collection and database design, described in detail in chapter \ref{chapter:data}. After dealing with data, we moved on prototyping our app to verify established requirements and gain insights how we want our app to look like. Further details about whole architecture and its implementation might be found in chapter \ref{chapter:system_design}.
    \par
    The final stage testing is documented in chapter \ref{chapter:testing}. This crucial phase allowed us to draw the final conclusions.